% !TeX root = ../main.tex

\section{การทดลองที่ 4: System Identification ของ Mass-Spring-Damper}
% \textit{ครอบคลุม Criteria: Part 1 System Modeling --- เขียน EOM,
% ออกแบบการทดลองหา $C$, อธิบายหลักการหา $C$, แสดงผล Model Fit (รวม 3 คะแนน)}

%% ─────────────────────────────────────────────────────────────────────────────
%% ─────────────────────────────────────────────────────────────────────────────
\subsection{บทนำ}

ในการสร้างแบบจำลองทางคณิตศาสตร์ (State-Space Model) ของระบบ Mass-Spring-Damper (MSD) เพื่อนำไปใช้ออกแบบตัวประมาณค่าสถานะ (State Estimator) นั้น ความแม่นยำของพารามิเตอร์ทางกายภาพคือหัวใจสำคัญ จากการทดลองที่ 1 ได้ประเมินค่า Spring Constant เบื้องต้นไว้ที่ $K = 29.11$~N/m อย่างไรก็ตาม แบบจำลองจะสมบูรณ์ได้จำเป็นต้องทราบพารามิเตอร์ที่ควบคุมอัตราการสูญเสียพลังงานของระบบด้วย นั่นคือ Damping Coefficient ($C$) 

การหาค่า $C$ มีความซับซ้อนกว่าการหาค่า $K$ ด้วยเหตุผลทางพลศาสตร์ดังต่อไปนี้:
\begin{itemize}
    \item \textbf{ข้อจำกัดในการวัดโดยตรง:} แตกต่างจากแรงสปริงที่ขึ้นอยู่กับระยะตำแหน่ง (Position) แรงหน่วงในระบบเป็นสัดส่วนโดยตรงกับความเร็ว (Velocity) ซึ่งวัดค่าแบบสถิต (Static) ไม่ได้
    \item \textbf{ความซับซ้อนของแรงเสียดทาน:} ในระบบกลไกจริง พารามิเตอร์ $C$ ไม่ได้เกิดจากตัวหน่วงอากาศเพียงอย่างเดียว แต่ยังรวมผลกระทบจากแรงเสียดทาน (Coulomb และ Viscous Friction) ที่แฝงอยู่ในระบบด้วย
\end{itemize}



ดังนั้น วิธีที่เหมาะสมที่สุดคือการทำ \textbf{System Identification} โดยการออกแบบการทดลองสภาวะการสั่นอย่างอิสระ (Free Vibration) คือการดึงมวลออกจากจุดสมดุล (Equilibrium) แล้วปล่อยโดยไม่มีแรงภายนอกมากระทำ ทำให้ $F(t) = 0$ สมการการเคลื่อนที่ (EOM) จึงลดรูปเหลือความสัมพันธ์ระหว่าง $m$, $K$, และ $C$ จากนั้นบันทึกการตอบสนองของระบบจริง (Measured Data) เพื่อนำไปวิเคราะห์ย้อนกลับ

\textbf{การปรับปรุงแนวทางการประเมินพารามิเตอร์ (Parameter Estimation Strategy):}
\begin{itemize}
    \item ในเบื้องต้น ทฤษฎีกำหนดให้ยึดค่า $m$ และ $K$ คงที่ เพื่อค้นหาค่า $C$ เพียงตัวเดียว
    \item อย่างไรก็ตาม ในการทดลองที่ 4 และ 5 ได้มีการปรับเปลี่ยนชุดทดลองโดยเพิ่ม \textbf{เสาไกด์ (Guide Rail)} แนบชิดกับฐานมวลเพื่อให้เซนเซอร์อ่านค่าได้เสถียรขึ้น การสัมผัสนี้ทำให้เกิดแรงต้านทานเพิ่มเติมและส่งผลให้แนวการรับแรงดิ่งของสปริงเปลี่ยนแปลงไปจากสภาวะอิสระของการทดลองที่ 1
    \item เพื่อป้องกันความคลาดเคลื่อนของแบบจำลอง (Model Mismatch) การทดลองนี้จึงใช้ MATLAB Parameter Estimator ในการ\textbf{ค้นหาค่า $C$ และปรับแก้ค่า $K$ ไปพร้อมกัน} โดยให้ Algorithm ทำการ Optimization จนกว่าผลการจำลอง (Simulation) จะสอดคล้องกับข้อมูลที่วัดได้จริงมากที่สุด (Best Fit)
\end{itemize}

พารามิเตอร์ $K$ ที่อัปเดตใหม่ และค่า $C$ (Effective Damping) ที่ครอบคลุมแรงต้านทานทั้งหมดในระบบ จะถูกนำไปใช้สร้างเมทริกซ์พลวัต (Discrete State Matrix, $A_d$) สำหรับ 2-State Kalman Filter ในการทดลองที่ 5 ต่อไป
%% ─────────────────────────────────────────────────────────────────────────────
\subsection{วัตถุประสงค์}

\begin{itemize}
    \item เพื่อ Derive และเขียน Equation of Motion (EOM) ของระบบ
          Mass-Spring-Damper จาก Newton's Second Law อย่างสมบูรณ์
    \item เพื่อออกแบบการทดลอง Free Vibration ที่เหมาะสม
          สำหรับการเก็บข้อมูล Displacement เพื่อนำไป Identify พารามิเตอร์
    \item เพื่อประมาณค่า Damping Coefficient $C$ โดยใช้
          MATLAB Parameter Estimator (Simulink Design Optimization)
    \item เพื่อตรวจสอบความถูกต้องของ Model โดยเปรียบเทียบผลการ Simulate
          กับข้อมูลจริงผ่าน Fit Metrics
\end{itemize}

%% ─────────────────────────────────────────────────────────────────────────────
\subsection{สมมติฐาน}

\begin{itemize}
    \item การตอบสนองของระบบภายใต้ Free Vibration จะมีลักษณะ Under-Damped Second-Order Response ($\zeta < 1$) เนื่องจากระบบ MSD ในชุดทดลองนี้มีการแกว่งที่สังเกตได้ชัดเจน
    \item Parameter Estimator จะค้นพบค่า $K$ ที่เปลี่ยนไปจากเดิมเล็กน้อย และค่า $C$ ที่ครอบคลุมทั้งแรงต้านอากาศและแรงเสียดทาน (Friction) จากเสาไกด์
    \item ค่าพารามิเตอร์ทั้งสองที่ประเมินได้ เมื่อประกอบกับมวล $m$ จะสามารถสร้างแบบจำลองที่จำลองพฤติกรรมของระบบจริงได้อย่างแม่นยำ
\end{itemize}

%% ─────────────────────────────────────────────────────────────────────────────
\subsection{ตัวแปร}

\begin{itemize}
    \item \textbf{ตัวแปรอิสระ:} Initial Displacement $x_0$ (m) คือจุดที่สปริงหดจากจุด Equilibrium (ประมาณ 7.2 cm)
    \item \textbf{ตัวแปรตาม:} Displacement ตามเวลา $x(t)$ วัดด้วย HC-SR04, \textbf{ค่า Spring Constant $K$ (N/m) และ Damping Coefficient $C$ (N$\cdot$s/m) ที่ประเมินได้}, Damping Ratio $\zeta$
    \item \textbf{ตัวแปรควบคุม:} มวล $m$ (kg) คงที่ตลอดการทดลอง, Sampling Time 10~ms, ความเร็วเริ่มต้น $\dot{x}(0) = 0$
\end{itemize}

%% ─────────────────────────────────────────────────────────────────────────────
\subsection{เอกสารและงานวิจัยที่เกี่ยวข้อง}

\subsubsection{Equation of Motion ของ Mass-Spring-Damper}

เมื่อพิจารณามวล $m$ ในระบบ Mass-Spring-Damper แบบแนวดิ่ง
\begin{figure}[htbp]
    \centering
    \includegraphics[width=0.3\textwidth]{exp4_fbd_msd.png}
    \caption{Free Body Diagram ของมวลในระบบ Mass-Spring-Damper
             แสดงแรงที่กระทำต่อมวลในทิศทาง $x$ (แนวดิ่ง)}
    \label{fig:fbd_msd}
\end{figure}    
มีแรงกระทำต่อมวล 3 แรงหลัก ได้แก่ แรงจากสปริง $-Kx$,
แรงหน่วง $-C\dot{x}$, และแรงภายนอก $F(t)$
เมื่อนำ Newton's Second Law ($\Sigma F = ma$) มาใช้จะได้:

\begin{equation}
    m\ddot{x} = -Kx - C\dot{x} + F(t)
    \label{eq:newton_msd}
\end{equation}

จัดรูปใหม่ได้ EOM มาตรฐานของระบบ Second-Order:

\begin{equation}
    m\ddot{x} + C\dot{x} + Kx = F(t)
    \label{eq:eom_msd}
\end{equation}

โดยที่ $m$ คือมวล (kg), $C$ คือ Damping Coefficient (N$\cdot$s/m),
$K$ คือ Spring Constant (N/m), $x$ คือการกระจัดจาก Equilibrium (m),
และ $F(t)$ คือแรงภายนอก (N) ซึ่งในกรณี Free Vibration มีค่า $F(t) = 0$

\subsubsection{เหตุผลที่ใช้ Displacement จาก Equilibrium เป็น State}

ในการสร้างแบบจำลองทางคณิตศาสตร์ การเลือกจุดอ้างอิง (Reference Point) มีผลอย่างมากต่อความซับซ้อนของสมการ State-Space เพื่อให้เข้าใจถึงประโยชน์ของการใช้ "การกระจัดจากจุดสมดุล" ($x$) แทน "ตำแหน่งสัมบูรณ์" ($x_\text{abs}$) สามารถพิจารณาได้จากลำดับดังต่อไปนี้



\textbf{1. นิยามของจุดอ้างอิงทางกายภาพ}
\begin{itemize}
    \item \textbf{$x_\text{abs}$ (Absolute Position):} ตำแหน่งที่วัดจากจุดที่สปริงยังไม่ยืดตัว (Unstretched Position)
    \item \textbf{$x_\text{eq}$ (Static Deflection):} ระยะยืดสถิต ซึ่งเป็นจุดที่มวลแขวนอยู่นิ่งๆ โดยที่แรงสปริงดึงกลับมีค่าเท่ากับน้ำหนักของมวลพอดี ($K x_\text{eq} = mg$)
    \item \textbf{$x$ (Displacement):} การกระจัดที่วัดจากจุดสมดุลใหม่ ซึ่งมีความสัมพันธ์คือ $x = x_\text{abs} - x_\text{eq}$
\end{itemize}

\textbf{2. การลดรูปสมการคณิตศาสตร์}

หากพิจารณา EOM โดยใช้ตำแหน่งสัมบูรณ์ $x_\text{abs}$ จะมีแรงโน้มถ่วง $mg$ ปรากฏเป็นค่าคงที่ในสมการ:
\begin{equation}
    m\ddot{x}_\text{abs} + C\dot{x}_\text{abs} + K x_\text{abs} = mg
    \label{eq:eom_absolute}
\end{equation}

เนื่องจากระยะ $x_\text{eq}$ เป็นค่าคงที่ การหาอนุพันธ์เทียบกับเวลาจึงมีค่าเท่ากับศูนย์ ส่งผลให้ความเร็วและความเร่งในทั้งสองระบบพิกัดมีค่าเท่ากัน ($\dot{x} = \dot{x}_\text{abs}$ และ $\ddot{x} = \ddot{x}_\text{abs}$) เมื่อแทนค่า $x_\text{abs} = x + x_\text{eq}$ ลงในสมการ~\eqref{eq:eom_absolute} จะได้:
\begin{equation}
    m\ddot{x} + C\dot{x} + K(x + x_\text{eq}) = mg
\end{equation}

กระจายเทอม $K$ และแทนค่าความสัมพันธ์สถิตยศาสตร์ $K x_\text{eq} = mg$:
\begin{equation}
    m\ddot{x} + C\dot{x} + Kx + mg = mg \quad \Rightarrow \quad m\ddot{x} + C\dot{x} + Kx = 0
    \label{eq:equilibrium_cancel}
\end{equation}

\textbf{3. ประโยชน์ในเชิงวิศวกรรมควบคุม (Control Theory)}

การที่เทอม $mg$ หายไปจากสมการ~\eqref{eq:equilibrium_cancel} มีนัยสำคัญอย่างมากต่อการออกแบบตัวควบคุมและตัวประมาณค่า:
\begin{itemize}
    \item \textbf{หลีกเลี่ยง Constant Disturbance:} หากใช้ $x_\text{abs}$ สมการ State-Space จะอยู่ในรูป $\dot{X} = AX + Bu + D$ (โดย $D$ คือเวกเตอร์ของแรงโน้มถ่วง) ซึ่งทำให้ระบบสูญเสียความเป็น Linear Homogeneous การแก้ปัญหานี้บังคับให้เราต้องถือว่าแรงโน้มถ่วงเป็น Input ($u$) ค่าคงที่ตลอดเวลา ซึ่งสร้างความซับซ้อนให้กับระบบเมทริกซ์โดยไม่จำเป็น
    \item \textbf{ลดความซับซ้อนของ Kalman Filter:} เมื่อย้ายจุดกำเนิดมาที่จุดสมดุล แรงโน้มถ่วงจะถูก "หักล้าง" ด้วยแรงตึงสปริงเริ่มต้น (Pre-load) ไปแล้ว ระบบจะเหลือเพียงพฤติกรรมทางพลศาสตร์ (Dynamic Behavior) ล้วนๆ ทำให้ State-Space Model อยู่ในรูปแบบ $\dot{X} = AX$ มาตรฐาน ส่งผลให้การเขียนโค้ดสำหรับ Prediction Step ใน Kalman Filter ตรงไปตรงมา และไม่ต้องเขียนอัลกอริทึมจัดการ Input Offset ชดเชยแรงโน้มถ่วงเพิ่มเติม
\end{itemize}
\subsubsection{Damping Ratio และประเภทของ Response}

Damping Ratio $\zeta$ เป็นพารามิเตอร์ไร้มิติที่บอกว่าระบบสูญเสียพลังงาน
และหยุดการแกว่งได้เร็วแค่ไหน นิยามโดยเปรียบเทียบ $C$ จริงกับ
Critical Damping Coefficient $C_c = 2\sqrt{Km}$:

\begin{equation}
    \zeta = \frac{C}{C_c} = \frac{C}{2\sqrt{Km}}
    \label{eq:damping_ratio}
\end{equation}

และ Natural Frequency $\omega_n$ นิยามเป็น:

\begin{equation}
    \omega_n = \sqrt{\frac{K}{m}}
    \label{eq:natural_freq}
\end{equation}

ค่า $\zeta$ กำหนดลักษณะของ Response โดยตรง:

\begin{figure}[htbp]
    \centering
    \includegraphics[width=0.7\textwidth]{exp4_damp_ratio.png}
    \caption{ลักษณะของการตอบสนองของระบบแบบ Step Input
             ตามค่า Damping Ratio $\zeta$}
    \label{fig:damp_ratio}
\end{figure}    

\begin{itemize}
    \item \textbf{Under-Damped ($\zeta < 1$):} ระบบแกว่งข้าม Equilibrium
          ไปมา โดยแอมพลิจูดลดลงแบบ Exponential จนหยุดนิ่ง
          รากของสมการลักษณะเฉพาะเป็นจำนวนเชิงซ้อน Conjugate
    \item \textbf{Critically-Damped ($\zeta = 1$):} ระบบกลับสู่ Equilibrium
          เร็วที่สุดโดยไม่มีการแกว่ง รากซ้ำกัน (Real and Repeated Roots)
    \item \textbf{Over-Damped ($\zeta > 1$):} ระบบไม่แกว่ง แต่กลับสู่
          Equilibrium ช้ากว่า Critically-Damped
          รากต่างกัน (Real and Distinct Roots)
\end{itemize}

\subsubsection{MATLAB Parameter Estimator (Simulink Design Optimization)}

MATLAB Parameter Estimator ทำงานโดยใช้ Algorithm ทางคณิตศาสตร์
(เช่น Nonlinear Least Squares) เพื่อค้นหาค่าพารามิเตอร์ที่ทำให้
ผลต่างระหว่าง Simulation กับข้อมูลจริงมีค่าน้อยที่สุด
โดยใช้ Objective Function:

\begin{equation}
    J(K,\,C) = \sum_{k=1}^{N}
               \left[ x_\text{measured}(t_k)
                    - x_\text{simulated}(t_k;\,K,\,C) \right]^2
    \label{eq:param_est_cost}
\end{equation}

แม้ว่าการทดลองที่ 1 จะได้ค่า $K = 29.11$ N/m มาแล้ว
แต่เนื่องจากการทดลองที่ 4 ใช้ Setup ที่ต่างออกไปโดยมีการเพิ่ม
Guide Rail เพื่อบังคับให้มวลเคลื่อนที่ในแนวดิ่งอย่างเคร่งครัด
การเปลี่ยนแปลงนี้ทำให้แนวการรับแรงของสปริงเปลี่ยนไป
ส่งผลให้ค่า $K$ ที่มีผลต่อพลศาสตร์ในแนวดิ่งเปลี่ยนแปลงตามด้วย
ดังนั้นจึงตัดสินใจปล่อยให้ Estimator ค้นหาค่า $K$ และ $C$ พร้อมกัน
เพื่อให้ได้พารามิเตอร์ที่สอดคล้องกับ Setup จริงของการทดลองนี้
มากที่สุด แทนที่จะ Fix $K$ ไว้ที่ค่าเดิมซึ่งวัดจาก Setup ที่ต่างกัน

Algorithm จะค้นหาคู่ $(K,\,C)$ ที่ทำให้ $J(K,C)$ มีค่าต่ำที่สุด
สำหรับการสร้าง Simulink Model ต้องจัดรูป EOM
ให้อยู่ในรูปที่แก้หาความเร่งได้:

\begin{equation}
    \ddot{x} = \frac{1}{m}\!\left(-C\dot{x} - Kx\right)
    \quad \text{(กรณี Free Vibration: }F(t) = 0\text{)}
    \label{eq:simulink_accel}
\end{equation}

จากนั้น Define State $v = \dot{x}$ เพื่อแปลงเป็น First-Order
State Equations สำหรับ Simulink:

\begin{equation}
    \dot{x} = v, \qquad
    \dot{v} = \frac{-Cv - Kx}{m}
    \label{eq:simulink_state}
\end{equation}

\subsection{การปรับปรุง Setup การทดลอง}
 
ในการทดลองที่ 1 สปริงถูกติดตั้งอย่างอิสระ ทำให้มวล (ฐานหนัก $m = 0.290$~kg)
มีโอกาสแกว่งออกนอกแนวดิ่งเล็กน้อย ส่งผลให้ HC-SR04
อ่านค่าระยะทางได้คลาดเคลื่อน เนื่องจากลำคลื่นอัลตราโซนิคต้องการพื้นผิวสะท้อน
ที่ตั้งฉากกับทิศทางการส่งคลื่น
 
การทดลองที่ 4 จึงปรับ Setup ใหม่โดยขยับเสาไกด์ (Guide Rail) เข้ามาชิดกับ
ฐานมวล ทำให้ฐานขนานกับพื้นและเคลื่อนที่ในแนวดิ่งได้ตรงยิ่งขึ้น
ส่งผลให้ Sensor อ่านค่าได้เสถียร อย่างไรก็ตาม การที่ฐานสัมผัสกับเสาไกด์
ทำให้เกิดแรงเสียดทาน (Friction) ในแนวดิ่งเพิ่มเติม ซึ่งส่งผลสองประการดังนี้
 
\begin{itemize}
    \item \textbf{ค่า $C$ ที่ Estimate ได้สูงกว่าความเป็นจริงของ Damper อย่างเดียว:}
          $C_\text{eff}$ ที่ได้จาก Parameter Estimator สะท้อนผลรวมของ
          Damping จริงของ Damper และ Friction ของ Guide Rail
          จึงถือเป็น Effective Damping Coefficient
          ที่ตรงกับระบบที่ใช้ในการทดลองที่ 5 มากที่สุด
 
    \item \textbf{ค่า $K$ ที่ Estimate ได้เปลี่ยนจาก Exp 1:}
          การปรับ Setup ทำให้ระยะ Pre-load ของสปริงและแนวการรับแรงเปลี่ยนไปบางส่วน
          ดังนั้น Parameter Estimator จึง Estimate ทั้ง $K$ และ $C$
          พร้อมกันจากข้อมูล Free Vibration จริงของ Setup นี้
          แทนที่จะ Fix $K = 29.11$~N/m จากการทดลองที่ 1
\end{itemize}
 
Simulink Model ที่สร้างขึ้นสอดคล้องกับ Setup นี้ โดย Output ของ Model
คือ Absolute Distance จาก Sensor (cm) ไม่ใช่ Displacement จาก Equilibrium
จึงต้องบวก Offset $x_\text{eq} = 11.3$~cm เข้าไปใน Model
ตามระยะที่ Sensor อ่านได้ที่ตำแหน่ง Equilibrium จริง

\subsubsection{การวิเคราะห์และคัดกรองข้อมูลผิดปกติ (Outlier Analysis)}

ในการทดลองทางวิศวกรรมและการระบุพารามิเตอร์ของระบบ (System Identification) ข้อมูลที่ได้จากการทดลองซ้ำ (Multiple Trials) มักมีการกระจายตัวอยู่ในช่วงระดับหนึ่งเนื่องจากสัญญาณรบกวน (Noise) และความไม่แน่นอนของระบบ อย่างไรก็ตาม ในบางกรณีอาจเกิด "ข้อมูลผิดปกติ" (Outlier) ซึ่งเป็นจุดข้อมูลที่มีค่าเบี่ยงเบนออกไปจากกลุ่มข้อมูลส่วนใหญ่อย่างสุดโต่ง สาเหตุอาจเกิดจากความผิดพลาดชั่วขณะของเซนเซอร์ การรบกวนจากสภาพแวดล้อม หรือพฤติกรรมทางกายภาพที่ผิดปกติ (เช่น การที่ฐานมวลเสียดสีกับเสาไกด์อย่างรุนแรงในบางรอบการทดลอง) 

หากนำข้อมูลที่เป็น Outlier มาร่วมคำนวณหาค่าเฉลี่ย (Mean) หรือความแปรปรวน (Variance) จะทำให้ค่าสถิติเหล่านั้นบิดเบือนไปจากความเป็นจริง (Skewed) ส่งผลให้แบบจำลองทางคณิตศาสตร์มีความคลาดเคลื่อน ดังนั้นจึงต้องมีกระบวนการทางสถิติที่น่าเชื่อถือในการตรวจจับและคัดกรองข้อมูลเหล่านี้ออกก่อนนำไปใช้งาน

\begin{figure}[htbp]
    \centering
    \includegraphics[width=0.7\textwidth]{bloxplot.png} 
    \caption{การวิเคราะห์ข้อมูลผิดปกติ (Outlier Analysis) ของค่า Damping Coefficient $C$ ด้วยวิธี Interquartile Range (IQR)}
    \label{fig:outlier_boxplot}
\end{figure}

\textbf{วิธีการพิสัยระหว่างควอร์ไทล์ (Interquartile Range: IQR)}
วิธี IQR เป็นเครื่องมือทางสถิติที่ได้รับความนิยมอย่างแพร่หลายในการระบุตำแหน่งของ Outlier เนื่องจากเป็นวิธีที่ทนทาน (Robust) ต่อข้อมูลที่มีความเบ้ (Skewness) และไม่จำเป็นต้องตั้งสมมติฐานว่าข้อมูลมีการแจกแจงแบบปกติ (Normal Distribution) โดยมีหลักการและขั้นตอนทางคณิตศาสตร์ดังนี้:

\begin{itemize}
    \item \textbf{การหาตำแหน่งควอร์ไทล์ (Quartiles):} นำชุดข้อมูลมาเรียงลำดับจากน้อยไปมาก จากนั้นแบ่งข้อมูลออกเป็น 4 ส่วนเท่าๆ กัน ค่าที่แบ่งส่วนที่ 1 และ 2 คือ ควอร์ไทล์ที่ 1 ($Q_1$ หรือ 25th Percentile) และค่าที่แบ่งส่วนที่ 3 และ 4 คือ ควอร์ไทล์ที่ 3 ($Q_3$ หรือ 75th Percentile)
    \item \textbf{การคำนวณระยะ IQR:} พิสัยระหว่างควอร์ไทล์คือระยะห่างระหว่าง $Q_3$ และ $Q_1$ ซึ่งครอบคลุมช่วงข้อมูลกึ่งกลาง 50\% ของข้อมูลทั้งหมด
    \begin{equation}
        \text{IQR} = Q_3 - Q_1
        \label{eq:iqr_calc}
    \end{equation}
    \item \textbf{การสร้างขอบเขต (Fences):} ใช้ระยะ IQR ในการกำหนดขอบเขตบนและล่างสำหรับคัดกรองข้อมูล โดยทั่วไปนิยมใช้ตัวคูณ 1.5 สำหรับข้อมูลผิดปกติระดับมาตรฐาน (Mild Outliers) 
    \begin{align}
        \text{Lower Fence} &= Q_1 - 1.5 \times \text{IQR} \label{eq:lower_fence} \\
        \text{Upper Fence} &= Q_3 + 1.5 \times \text{IQR} \label{eq:upper_fence}
    \end{align}
    \item \textbf{เกณฑ์การตัดสิน:} ข้อมูลใดๆ ที่มีค่าน้อยกว่า Lower Fence หรือมีค่ามากกว่า Upper Fence จะถูกพิจารณาว่าเป็น Outlier ทางสถิติ และมีความชอบธรรมทางวิชาการที่จะถูกแยกออกจากการคำนวณพารามิเตอร์ของระบบ
\end{itemize}

%% ─────────────────────────────────────────────────────────────────────────────
\subsection{ขั้นตอนการทดลอง}

\textbf{ส่วนที่~1: เตรียมชุดทดลองและเก็บข้อมูล Free Vibration}

\begin{enumerate}
    \item ประกอบชุด Mass-Spring-Damper \textbf{โดยปรับเสาไกด์ให้แนบกับฐานมวลพอดี} วาง HC-SR04 ด้านล่างชี้ขึ้นเพื่อวัด Displacement ในแนวดิ่ง
    \item ชั่งมวล $m$ ด้วยเครื่องชั่ง บันทึกค่าเป็น kg
    \item ตั้งระบบให้มวลอยู่ที่ Equilibrium Position ยืนยันว่า HC-SR04 อ่านค่า Baseline คงที่ กำหนดตำแหน่งนี้เป็น $x = 0$
    \item ดึงมวลขึ้นจาก Equilibrium เป็นระยะประมาณ 7.2 cm จนสปริงหดแล้วปล่อยโดยไม่ออกแรง ($\dot{x}(0) = 0$) บันทึก Oscillation จนกว่าจะ Damp Out ส่งข้อมูลผ่าน Serial 115200 baud
    \item ทำซ้ำการ Release 5 ครั้งเพื่อตรวจสอบ Repeatability บันทึกทุก Dataset แยกไฟล์
\end{enumerate}

\textbf{ส่วนที่~2: System Identification ด้วย MATLAB Parameter Estimator}

\begin{enumerate}
    \item Import ข้อมูล Displacement จาก Serial ลงใน MATLAB แปลงหน่วยเป็นเมตร
    \item สร้าง Simulink Model ของ MSD โดยใช้ State Equations: Input $u = 0$, Initial Condition $x(0) = x_0$, $v(0) = 0$
    \item เปิด Simulink Design Optimization \textbf{กำหนดให้ $K$ และ $C$ เป็น Parameter ที่ต้องการ Estimate พร้อมกัน} โดยตั้งค่าเริ่มต้น (Initial Guess) ของ $K = 29.11$~N/m (จาก Exp 1) และ Fix ค่า $m$ ไว้
    \item รัน Estimator จนกว่าจะ Converge บันทึกค่า $K$ และ $C$ ที่ได้
    \item คำนวณ $\zeta$ จากสมการ $\zeta = C / (2\sqrt{Km})$ และจัดประเภท Damping
    \item Overlay กราฟ Model Simulation กับ Measured Data บันทึกค่า RMSE และทำ Outlier Analysis
\end{enumerate}

\subsection{ผลการทดลอง}
 
\subsubsection{Equation of Motion}
 
จาก Newton's Second Law สำหรับระบบ Mass-Spring-Damper
ที่ใช้ Displacement จาก Equilibrium เป็น State:
 
\begin{equation}
    m\ddot{x} + C\dot{x} + Kx = 0
    \quad \text{(Free Vibration)}
    \label{eq:eom_free_vib}
\end{equation}
 
% โดยพารามิเตอร์ที่ทราบจากการทดลองที่ 1: $K_\text{Exp1} = 29.11$~N/m
% และมวล $m = 0.290$~kg
\subsubsection{Outlier Analysis ของผล Parameter Estimation}
 
จากการทำซ้ำ 5 Trials พบว่าค่า $C$ ของ Trial 8 มีค่าสูงผิดปกติ
จึงทำการตรวจสอบด้วยวิธี IQR (Interquartile Range)
 
\begin{table}[htbp]
    \centering
    \caption{ผลการ Estimate พารามิเตอร์ $K$ และ $C$ จาก 5 Trials
             (MATLAB Parameter Estimator, Simulink Design Optimization)}
    \label{tab:param_est_raw}
    \begin{tabular}{cccc}
        \hline
        \textbf{Trial} & \textbf{$K$ (N/m)} & \textbf{$C$ (N$\cdot$s/m)} & \textbf{หมายเหตุ} \\
        \hline
        1 & 33.5482 & 0.3961 & --- \\
        4 & 31.4963 & 0.3760 & --- \\
        5 & 33.2512 & 0.4184 & --- \\
        7 & 29.8704 & 0.3623 & --- \\
        8 & 30.9629 & 0.9792 & Outlier (IQR) \\
        \hline
    \end{tabular}
\end{table}
 
การตรวจสอบ Outlier ด้วยวิธี IQR สำหรับค่า $C$
มีขั้นตอนดังนี้
 
\begin{enumerate}
    \item เรียงค่า $C$ จากน้อยไปมาก:
          $\{0.3623,\; 0.3760,\; 0.3961,\; 0.4184,\; 0.9792\}$
    \item คำนวณ Quartile:
          $Q_1 = 0.3760$, $Q_3 = 0.4184$,
          $\text{IQR} = Q_3 - Q_1 = 0.0424$
    \item กำหนด Fence:
          \begin{align}
              \text{Lower fence} &= Q_1 - 1.5 \times \text{IQR}
                                  = 0.3760 - 0.0635 = 0.3125 \notag \\
              \text{Upper fence} &= Q_3 + 1.5 \times \text{IQR}
                                  = 0.4184 + 0.0635 = 0.4819 \notag
          \end{align}
    \item Trial 8 มีค่า $C = 0.9792 > 0.4819$ จึงถือเป็น Outlier
          และตัดออกจากการคำนวณค่าเฉลี่ย
\end{enumerate}
 
\subsubsection{สถิติของพารามิเตอร์ที่ได้}
 
\begin{table}[htbp]
    \centering
    \caption{สถิติของ $K$ และ $C$ ก่อนและหลังตัด Outlier}
    \label{tab:param_stats}
    \begin{tabular}{lccccc}
        \hline
        \textbf{กลุ่ม} & \textbf{$n$}
            & \textbf{$\bar{K}$ (N/m)} & \textbf{S.D.($K$)}
            & \textbf{$\bar{C}$ (N$\cdot$s/m)} & \textbf{S.D.($C$)} \\
        \hline
        ทั้งหมด (5 Trials) & 5 & 31.83 & 1.56 & 0.5064 & 0.2652 \\
        ไม่รวม Trial 8     & 4 & 32.04 & 1.71 & 0.3882 & 0.0245 \\
        \hline
    \end{tabular}
\end{table}
 
การตัด Outlier ทำให้ S.D. ของ $C$ ลดลงจาก 0.2652 เป็น 0.0245~N$\cdot$s/m
(ลดลง 90.8\%) ขณะที่ S.D. ของ $K$ ไม่เปลี่ยนแปลงมากนัก
ยืนยันว่า Trial 8 มีผลกระทบต่อ $C$ อย่างมีนัยสำคัญ
 
\subsubsection{Simulink Model และ Parameter Estimator}
 
\begin{figure}[htbp]
    \centering
    % TODO: แทนด้วยรูป Simulink Model จริง
    \includegraphics[width=0.90\textwidth]{exp4_simulink.png}
    \caption{Simulink Model ของระบบ MSD สำหรับ Parameter Estimation
             แสดงโครงสร้าง Feedback Loop ของ $\ddot{x} = (-C\dot{x} - Kx)/m$
             พร้อม Offset $x_\text{eq} = 11.3$~cm เพื่อแปลง Displacement
             เป็น Absolute Distance ที่ Sensor วัดได้}
    \label{fig:exp4_simulink}
\end{figure}
 
Model ถูกสร้างตามสมการ~\eqref{eq:simulink_state} โดยมี Offset
$x_\text{eq} = 11.3$~cm บวกเข้าไปที่ Output เนื่องจาก Sensor
อ่านค่า Equilibrium Position ได้ 11.3~cm จาก Sensor ถึงมวล
Parameter Estimator จึง Fit กับ Absolute Distance
ไม่ใช่ Displacement จาก Equilibrium โดยตรง
 
\begin{figure}[htbp]
    \centering
    % TODO: แทนด้วยกราฟ Repeatability จริง (3 Trials ซ้อนกัน)
    \includegraphics[width=0.85\textwidth]{exp4_repeatability.png}
    \caption{Time-Domain Plot ของ Displacement $x(t)$
             (แกน $x$: เวลา (s), แกน $y$: ระยะทาง (cm))}
    \label{fig:exp4_repeatability}
\end{figure}

\begin{table}[htbp]
\centering
\caption{ผลการประมาณค่าพารามิเตอร์และ Goodness-of-Fit จาก MATLAB Parameter Estimator}
\label{tab:param_fit}
\begin{tabular}{lcccccc}
\toprule
Trial & $K$ (N/m) & $C$ (N$\cdot$s/m) & RMSE (cm) & NRMSE & Fit\% & หมายเหตุ \\
\midrule
1 & 33.5482 & 0.3961 & 0.522 & 0.0585 & 48.8 & — \\
4 & 31.4963 & 0.3760 & 0.387 & 0.0399 & 62.4 & — \\
5 & 33.2512 & 0.4184 & 0.462 & 0.0524 & 54.1 & — \\
7 & 29.8704 & 0.3623 & 0.515 & 0.0451 & 56.7 & — \\
8 & 30.9629 & 0.9792 & 0.481 & 0.0475 & —    & Outlier (IQR) \\
\midrule
\multicolumn{2}{l}{เฉลี่ย (Trial 1,4,5,7)} & & 0.472 & 0.0490 & 55.5 & \\
\bottomrule
\end{tabular}
\end{table}

\begin{figure}[htbp]
    \centering
    \includegraphics[width=0.85\textwidth]{exp4_model_fit.png}
    \caption{Model-Fit Plot: Simulation จาก Parameter Estimator (เส้นประสีแดง)
             เทียบกับ Measured Data (เส้นทึบสีน้ำเงิน)
             (แกน $x$: เวลา (s), แกน $y$: ระยะทาง (cm))}
    \label{fig:exp4_model_fit}
\end{figure}

\subsubsection{พารามิเตอร์สุดท้ายที่ได้จากการทดลอง}

ใช้ค่าเฉลี่ยจาก 4 Trials ที่ผ่าน Outlier Test เป็นค่าสุดท้าย
สำหรับนำไปใช้ในการทดลองที่ 5

\begin{table}[htbp]
    \centering
    \caption{พารามิเตอร์สุดท้ายของระบบ Mass-Spring-Damper
             (เฉลี่ยจาก 4 Trials, ตัด Trial 8 ออกเป็น Outlier)}
    \label{tab:msd_params}
    \begin{tabular}{lcccc}
        \hline
        \textbf{พารามิเตอร์} & \textbf{สัญลักษณ์}
            & \textbf{หน่วย} & \textbf{ค่า} & \textbf{S.D.} \\
        \hline
        มวล                  & $m$          & kg           & 0.290    & --- \\
        Spring Constant      & $K$          & N/m          & 32.04    & 1.71 \\
        Damping Coefficient  & $C$          & N$\cdot$s/m  & 0.3882   & 0.0245 \\
        Natural Frequency    & $\omega_n$   & rad/s        & 10.51    & --- \\
        Natural Frequency    & $f_n$        & Hz           & 1.673    & --- \\
        Critical Damping     & $C_c$        & N$\cdot$s/m  & 6.097    & --- \\
        Damping Ratio        & $\zeta$      & ---          & 0.0637   & --- \\
        ประเภท Damping       & ---          & ---          & \multicolumn{2}{l}{Under-Damped} \\
        Fit Percentage       & Fit\%        & \%           & 55.5     & --- \\
        RMSE                 & ---          & cm           & 0.472    & --- \\
        \hline
    \end{tabular}
\end{table}

โดยที่

\begin{align}
    \omega_n &= \sqrt{\frac{K}{m}} = \sqrt{\frac{32.04}{0.290}}
               = 10.51 \text{ rad/s}
               \label{eq:omega_n_calc} \\[4pt]
    C_c      &= 2\sqrt{Km} = 2\sqrt{32.04 \times 0.290}
               = 6.097 \text{ N$\cdot$s/m}
               \label{eq:cc_calc} \\[4pt]
    \zeta    &= \frac{C}{C_c} = \frac{0.3882}{6.097}
               = 0.0637
               \label{eq:zeta_calc}
\end{align}

ค่า $\zeta = 0.0637 \ll 1$ ยืนยันว่าระบบอยู่ในสภาวะ Under-Damped
สอดคล้องกับที่สังเกตเห็นการแกว่งชัดเจนในการทดลอง

\subsubsection{การวิเคราะห์ความถูกต้องของแบบจำลอง}

จากตารางที่~\ref{tab:param_fit} ค่า RMSE เฉลี่ยของ 4 Trial ที่ผ่าน
Outlier Test อยู่ที่ 0.472~cm และค่า NRMSE อยู่ที่ 0.049
ซึ่งหมายความว่าความคลาดเคลื่อนเฉลี่ยของ Model อยู่ที่ประมาณ 4.9\%
ของพิสัยของสัญญาณ ค่า Fit\% เฉลี่ยอยู่ที่ 55.5\% โดย Trial 4
ให้ค่าสูงสุดที่ 62.4\% และ Trial 1 ให้ค่าต่ำสุดที่ 48.8\%
ซึ่งสะท้อนความแตกต่างของ Initial Displacement และความสม่ำเสมอ
ของการ Release ในแต่ละ Trial

ความคลาดเคลื่อนที่เกิดขึ้นมีสาเหตุหลักสามประการ ประการแรกคือ
สัญญาณรบกวนของ HC-SR04 ที่มี Variance สูงถึง 0.033~cm$^2$
ที่ระยะ 20~cm ทำให้ Measured Signal มี Noise ปนอยู่ตลอด
การ Fit กับ Smooth ODE Solution จึงมี Residual จาก Noise
โดยธรรมชาติ ประการที่สองคือ Friction จาก Guide Rail
ทำให้ Effective Damping ไม่คงที่ตามความเร็ว
Linear Model จึงไม่สามารถ Fit พฤติกรรมจริงได้อย่างสมบูรณ์
โดยเฉพาะ Trial 8 ที่มีค่า $C = 0.9792$~N$\cdot$s/m
ซึ่งสูงกว่า Trial อื่นมากกว่า 2 เท่า และถูกระบุเป็น Outlier
ด้วยวิธี IQR ประการที่สามคือ Linear Spring Model
สมมติว่า $K$ คงที่ตลอดช่วง Displacement
ซึ่งอาจไม่สมบูรณ์เมื่อสปริงทำงานนอกช่วง Linear Elastic

อย่างไรก็ตาม ค่า $\bar{K} = 32.04$~N/m และ
$\bar{C} = 0.3882$~N$\cdot$s/m ที่ได้มีความสม่ำเสมอเพียงพอ
(S.D. ของ $C$ ลดลง 90.8\% หลังตัด Outlier)
สำหรับการนำไปสร้าง Discrete State Matrix $A_d$
ใน 2-State Kalman Filter ของการทดลองที่~5

\subsection{สรุปผลการทดลอง}

\begin{itemize}
    \item ได้ค่า Spring Constant $\bar{K} = 32.04 \pm 1.71$~N/m
          และ Damping Coefficient $\bar{C} = 0.3882 \pm 0.0245$~N$\cdot$s/m
          จากการ Estimate 4 Trials ที่ผ่าน Outlier Test (ตัด Trial 8 ออก)

    \item ค่า Fit\% เฉลี่ยอยู่ที่ 55.5\% และ RMSE เฉลี่ยอยู่ที่ 0.472~cm
          โดย Trial 4 ให้ Fit\% สูงสุดที่ 62.4\% สะท้อนว่า Linear MSD Model
          สามารถอธิบายพฤติกรรมหลักของระบบได้ในระดับหนึ่ง
          แต่ยังมีความคลาดเคลื่อนจาก Nonlinear Friction และ Sensor Noise

    \item Trial 8 ถูกระบุเป็น Outlier ด้วยวิธี IQR โดยมีค่า
          $C = 0.9792$~N$\cdot$s/m ซึ่งเกิน Upper Fence ที่ 0.4819~N$\cdot$s/m
          อาจเกิดจากการสัมผัสระหว่างฐานมวลและเสาไกด์ที่รุนแรงผิดปกติ
          ในรอบนั้น

    \item ค่า $K$ ที่ Estimate ได้ (32.04~N/m) สูงกว่าค่าจากการทดลองที่ 1
          (29.11~N/m) ประมาณ 10\% เนื่องจาก Setup ที่เปลี่ยนแปลง
          จึงใช้ค่าใหม่นี้ใน State-Space Model ของการทดลองที่ 5

    \item ระบบอยู่ในสภาวะ Under-Damped ($\zeta = 0.0637$)
          สอดคล้องกับสมมติฐานและการสังเกตการแกว่งในการทดลอง

    \item ค่าพารามิเตอร์ $m = 0.290$~kg, $K = 32.04$~N/m,
          $C = 0.3882$~N$\cdot$s/m จะถูกนำไปสร้าง Discrete State Matrix $A_d$
          ใน 2-State Kalman Filter ของการทดลองที่ 5
\end{itemize}

\subsection{อภิปรายผล}
\begin{itemize}
    \item \textbf{ผลกระทบของ Friction ต่อค่า $C$ ที่ Estimate ได้:}
          ค่า $C_\text{eff} = 0.3882$~N$\cdot$s/m ที่ได้สะท้อนผลรวมของ
          Damping จาก Damper และ Friction จาก Guide Rail
          ซึ่งเป็น Effective Damping ของระบบจริงในการทดลองที่ 5
          การพยายามแยก Damping ออกจาก Friction ต้องการการทดลองเพิ่มเติม
          เช่น การวัด Friction Force โดยตรง ซึ่งอยู่นอกขอบเขตการทดลองนี้

    \item \textbf{ความแตกต่างของ $K$ ระหว่าง Exp 1 และ Exp 4:}
          ค่า $K$ เพิ่มขึ้นจาก 29.11 เป็น 32.04~N/m (เพิ่มขึ้น 10.1\%)
          เมื่อ Setup เปลี่ยน เนื่องจากการปรับตำแหน่งเสาไกด์
          ทำให้สปริงรับแรงในแนวดิ่งมากขึ้น ลด Component
          ของแรงในแนวเฉียงที่ไม่ได้ร่วมกับ Hooke's Law ในทิศดิ่ง
          ดังนั้น $K = 32.04$~N/m จึงแม่นยำกว่าสำหรับ Setup ของ Exp 4--5

    \item \textbf{หลักการทำงานของ Parameter Estimator:}
          Simulink Design Optimization ใช้ Algorithm Nonlinear Least Squares
          ปรับค่า $K$ และ $C$ ซ้ำๆ จนกว่า $J(K,C)$
          ตามสมการ~\eqref{eq:param_est_cost} จะมีค่าต่ำที่สุด
          กระบวนการนี้เทียบเท่าการค้นหาจุด Minimum ของ Error Surface
          ในปริภูมิ 2 มิติ ($K$--$C$)

    \item \textbf{เหตุผลที่ EOM ใช้ Equilibrium เป็น $x = 0$:}
          ตามสมการ~\eqref{eq:equilibrium_cancel} เมื่อใช้ตำแหน่งสัมบูรณ์
          EOM จะมีเทอม $mg$ ที่ไม่ Cancel ออกไป
          แต่เมื่อใช้การกระจัดจาก Equilibrium เทอม $mg$
          Balance กับ Spring Pre-load พอดี ทำให้ $u = 0$
          ใน Free Vibration อย่างไรก็ตาม เนื่องจาก Simulink Model
          ต้อง Fit กับ Absolute Distance จาก Sensor
          จึงต้องบวก Offset $x_\text{eq} = 11.3$~cm
          ที่ Output ของ Model

    \item \textbf{แหล่งที่มาของความคลาดเคลื่อนของ Model:}
          \begin{itemize}
              \item Friction ระหว่างฐานมวล
                    (velocity-dependent และ position-dependent)
                    ทำให้ Linear Damping Model มีความคลาดเคลื่อน
                    โดยเฉพาะใน Trial 8 ที่มีค่า $C$ สูงผิดปกติ
              \item ความไม่เชิงเส้นของสปริงที่ Displacement ใหญ่
              \item Noise ของ HC-SR04 ใน Training Data
                    ส่งผลต่อ Cost Function ของ Estimator
          \end{itemize}

    \item \textbf{ผลกระทบต่อ 2-State KF ในการทดลองที่ 5:}
          ค่า $K$ และ $C$ ที่ได้มีผลโดยตรงต่อ Discrete State Matrix $A_d$
          \[
              A_d = \begin{bmatrix} 1 & T_s \\ -\frac{K}{m}T_s & 1-\frac{C}{m}T_s \end{bmatrix}
          \]
          หากค่าเหล่านี้มีความคลาดเคลื่อน Prediction Step ของ KF
          จะ Propagate ความคลาดเคลื่อนสะสม ทำให้ต้องเพิ่ม $Q$
          เพื่อชดเชย Model Uncertainty
\end{itemize}